
% !TeX root = main.tex
\newpage
\clearpage
\begin{center}
\Large\textbf{Conventions Section\\ (included only in rough draft. For documenting standards, outline, and rules)}
\end{center}
\section{Manuscript Outline}

\begin{verbatim}
Abstract
What was done?
Why is it important?
What was achieved?

1. Introduction
	Why are particle systems important?
		Paragraphs 1–2

	What are particle systems used for?
		Paragraphs 1–2

	Why are large particle systems computationally limited?
		Paragraphs 3–5

	What is the contribution of gPCD?
		Paragraphs 6–7

	Paragraph 1

	   Purpose:
	   - Define particle methods.
	   - Explain why particle methods are useful.
	   - Introduce emergent macroscopic behavior.

	   Contains:
	   - Particle methods model discrete interacting entities.
	   - Applicable to solids, fluids, granular materials,
		 multiphase flows, molecular systems, and other applications.
	   - Individual entities are represented explicitly.
	   - Macroscopic behavior emerges from local interactions.
	   - Particle methods are well suited for studying collective behavior.

	   Presentation Bullets:
	   - Discrete interacting particles.
	   - Local interactions create emergent behavior.
	   - Applicable to many scientific and engineering systems.

	Paragraph 2

	   Purpose:
	   - Explain why large particle simulations are important.
	   - Motivate particle simulation as a discovery tool.
	   - Connect simulation to scientific advancement and AI.

	   Contains:
	   - Particle simulations support scientific discovery.
	   - New physical behavior may only emerge in large-scale simulations.
	   - Machine-learning methods require representative training data.
	   - Simulations generate training datasets.
	   - Simulations enable investigation of poorly understood phenomena.
	   - Simulations complement experiments and analysis.

	   Presentation Bullets:
	   - Large-scale simulations enable scientific discovery.
	   - Emergent phenomena may require many interacting particles.
	   - Simulations generate datasets for machine learning.
	   - Simulations help investigate previously unobserved behavior.

	Paragraph 3

	   Purpose:
	   - Explain why large particle counts are desirable.
	   - Explain why large particle simulations become expensive.
	   - Introduce the need for efficient particle-processing methods.

	   Contains:
	   - Simulation fidelity increases with particle count.
	   - Realistic simulations may require millions or billions of particles.
	   - Larger particle counts improve spatial resolution.
	   - Larger particle counts improve geometric representation.
	   - Larger particle counts enable multiscale investigations.
	   - Increased particle counts substantially increase computational cost.
	   - Efficient particle-processing algorithms are required for large-scale simulation.

	   Presentation Bullets:
	   - High-fidelity simulations require millions to billions of particles.
	   - Larger particle counts improve resolution and physical realism.
	   - Computational cost grows rapidly with problem size.
	   - Efficient particle-processing algorithms are essential.
   
	Paragraph 4

		Purpose:
		- Identify collision detection as a major computational bottleneck.
		- Explain why naïve collision detection is impractical.
		- Introduce broad-phase collision detection.
		- Motivate the need for efficient parallel broad-phase methods.

		Contains:
		- Collision detection is a primary limitation of large-scale particle simulation.
		- Naïve collision detection scales as O(N²).
		- Direct all-to-all evaluation is impractical for large particle counts.
		- Broad-phase methods reduce collision-detection cost.
		- Broad-phase methods restrict testing to nearby particles.
		- Efficient broad-phase construction remains challenging on massively parallel hardware.

		Presentation Bullets:
		- Collision detection is a major performance bottleneck.
		- Naïve collision detection scales as O(N²).
		- Broad-phase methods reduce candidate pairs.
		- Efficient parallel broad-phase construction remains challenging.
	   
 Paragraph 5

	Purpose:
		- Introduce GPUs as a platform for large-scale particle simulation.
		- Explain that hardware capability alone does not solve the problem.
		- Identify occupancy generation as the remaining bottleneck.
		- Transition toward the architectural motivation for gPCD.

		Contains:
		- GPUs provide high computational throughput.
		- GPUs provide high memory bandwidth.
		- GPUs are well suited for particle simulation.
		- Spatial occupancy structures remain difficult to construct efficiently.
		- Occupancy generation remains a bottleneck at large scales.

		Presentation Bullets:
		- GPUs provide massive parallelism.
		- GPUs offer high memory bandwidth.
		- Particle simulation maps naturally to GPU hardware.
		- Occupancy generation remains a major bottleneck.  
   
	Paragraph 6

		Purpose:
		- Introduce gPCD.
		- State the primary contribution of the work.
		- Summarize the key architectural ideas.
		- Explain the significance of the framework.

		Contains:
		- gPCD is a GPU-driven particle collision-detection framework.
		- gPCD uses the graphics pipeline for occupancy generation.
		- gPCD constructs a uniform cell-array occupancy structure.
		- Occupancy generation is organized around particle corner locations.
		- GPU rasterization hardware is used for broad-phase construction.
		- The framework maintains a simple and scalable data structure.
		- The framework supports large particle counts on a single CPU–GPU system.
		- The framework provides a foundation for future large-scale particle simulations.

		Presentation Bullets:
		- GPU-driven particle collision detection (gPCD).
		- Graphics-pipeline occupancy generation.
		- Corner-based occupancy construction.
		- Uniform cell-array broad phase.
		- Scalable single CPU–GPU execution.
	
	Paragraph 7

		Purpose:
		- Clearly define the novelty of the work.
		- Distinguish gPCD from existing occupancy structures.
		- Explain the architectural contribution.
		- Explain why the graphics pipeline is used.

		Contains:
		- The contribution is an execution architecture rather than a new data structure.
		- gPCD integrates occupancy generation directly into the graphics pipeline.
		- Rendering and occupancy construction occur concurrently.
		- Occupancy generation produces the potentially colliding set (PCS).
		- Graphics-pipeline resources participate in collision detection.
		- Resources traditionally reserved for rendering are repurposed for occupancy generation.

		Presentation Bullets:
		- Novelty: execution architecture, not a new data structure.
		- Graphics-pipeline occupancy generation.
		- Concurrent rendering and occupancy construction.
		- PCS generation integrated into rendering workflow.
		- Repurposes graphics resources for collision detection.
	
	
	
	
2. Background and Related Work
   What must the reader know first?

   Section Introduction

	   Purpose:
	   - Introduce the scope of the background section.
	   - Explain that collision-detection methods differ in both
		 data structures and execution architectures.
	   - Establish the relationship between collision detection
		 and rendering.
	   - Prepare the reader for the architectural gap discussion.

	   Contains:
	   - Many collision-detection approaches have been proposed.
	   - Approaches differ in both data structures and execution architectures.
	   - Collision detection may be organized differently relative to rendering.
	   - The section reviews methodologies, architectures,
		 and rendering strategies.
	   - The section identifies the architectural gap addressed by gPCD.

	   Presentation Bullets:
	   - Collision-detection methods vary in structure and execution model.
	   - Rendering and collision detection are often separated.
	   - Existing architectures leave an unexploited opportunity.
	   - Background motivates the architectural contribution of gPCD.
   
	2.1 Collision-Detection Frameworks
		What is occupancy generation and why is it needed?

	Paragraph 1

	   Purpose:
	   - Introduce the broad-phase/narrow-phase organization.
	   - Define the two major stages of collision detection.
	   - Establish where gPCD contributes.
	   - Explain why broad-phase generation is important.

	   Contains:
	   - Collision detection is commonly divided into broad-phase and narrow-phase stages.
	   - Broad phase generates collision candidates.
	   - Narrow phase evaluates candidate collisions.
	   - This organization is widely used across collision-detection frameworks.
	   - gPCD's primary contribution occurs in broad-phase occupancy generation.
	   - Narrow phase completes the overall collision-detection process.

	   Presentation Bullets:
	   - Collision detection is typically divided into two stages.
	   - Broad phase generates collision candidates.
	   - Narrow phase evaluates candidate collisions.
	   - gPCD focuses on broad-phase occupancy generation.
	   

	Paragraph 2

		Purpose:
		- Review common GPU broad-phase methodologies.
		- Show that many approaches exist for candidate generation.
		- Explain the common objective of broad-phase methods.
		- Transition from collision-detection methods to execution architectures.

		Contains:
		- Numerous GPU broad-phase methods have been proposed.
		- Common approaches include:
		  - Uniform grids
		  - Spatial hashing
		  - Hierarchical spatial subdivision
		  - Sweep and prune
		  - Image-space methods
		  - Other acceleration techniques
		- Methods differ in how spatial information is organized and queried.
		- All broad-phase methods seek to reduce narrow-phase evaluations.
		- Research has expanded from algorithms to execution architectures.
		- GPU computing has increased interest in architectural design.

		Presentation Bullets:
		- Many GPU broad-phase methodologies exist.
		- Methods differ in spatial organization and query strategy.
		- All seek to reduce narrow-phase collision tests.
		- Recent work increasingly focuses on execution architectures.

	2.2 Collision-Detection Architectures
		How have rendering and collision detection traditionally
		been organized?

   Paragraph 1

	   Purpose:
	   - Introduce execution architectures as distinct from collision-detection methods.
	   - Describe the historical evolution of collision-detection execution.
	   - Explain the migration from CPU to GPU implementations.
	   - Establish the architectural spectrum considered in the literature.

	   Contains:
	   - Collision-detection methodologies can be implemented using different execution architectures.
	   - Early systems performed broad-phase and narrow-phase operations on the CPU.
	   - Collision detection and rendering were traditionally separate processes.
	   - Increasing particle counts motivated GPU acceleration.
	   - Programmable GPUs enabled portions of the collision-detection pipeline to migrate to the GPU.
	   - GPU parallelism became a major driver of architectural evolution.
	   - Existing architectures span:
		   - CPU-resident implementations
		   - GPU-assisted implementations
		   - Fully GPU-resident implementations

	   Presentation Bullets:
	   - Execution architecture is distinct from collision-detection methodology.
	   - Early systems were CPU-based.
	   - Collision detection and rendering were traditionally independent.
	   - GPU programmability enabled migration of collision detection to the GPU.
	   - Architectures now range from CPU-resident to fully GPU-resident.

	Paragraph 2

		Purpose:
		- Define the major architectural categories.
		- Explain how work is distributed between CPU and GPU.
		- Describe the progression toward increased GPU utilization.
		- Introduce data movement as an architectural concern.

		Contains:
		- CPU-resident architectures perform broad-phase and narrow-phase operations on the CPU.
		- Simulation state is transferred to the rendering system after collision processing.
		- GPU-assisted architectures offload selected collision-detection tasks to the GPU.
		- CPU-managed execution and data structures remain part of the workflow.
		- GPU-resident architectures perform broad-phase and narrow-phase operations entirely on the GPU.
		- GPU-resident approaches reduce host-device communication.
		- GPU-resident approaches improve parallel utilization.
		- GPU-resident approaches reduce data movement overhead.

		Presentation Bullets:
		- CPU-resident: collision detection performed entirely on the CPU.
		- GPU-assisted: selected operations offloaded to the GPU.
		- GPU-resident: broad phase and narrow phase executed on the GPU.
		- Increasing GPU utilization reduces data-transfer overhead.


	Paragraph 3

		Purpose:
		- Identify a common limitation across existing architectures.
		- Explain that rendering and collision detection are usually separate.
		- Highlight the traditional role of the rendering pipeline.
		- Motivate examination of rendering-collision integration.

		Contains:
		- Existing architectures differ in execution location and data movement.
		- Collision detection and rendering are typically separate processes.
		- Collision-detection work is performed in dedicated computational stages.
		- The rendering pipeline primarily consumes simulation state.
		- Rendering generally does not participate in collision detection.
		- This separation motivates investigation of rendering-collision interaction.

		Presentation Bullets:
		- Most architectures separate rendering and collision detection.
		- Collision detection uses dedicated compute stages.
		- Rendering typically consumes simulation results.
		- Rendering rarely contributes to collision detection.
		- This separation motivates the gPCD architecture.   
   
	2.3 Rendering and Collision Detection

	Paragraph 1

		Purpose:
		- Define the traditional roles of collision detection and rendering.
		- Explain why the two systems have historically been separated.
		- Establish the conventional simulation-to-rendering workflow.
		- Prepare the reader for alternative rendering-based architectures.

		Contains:
		- Collision detection identifies particle interactions.
		- Collision detection advances simulation state.
		- Rendering visualizes simulation results.
		- Collision detection and rendering perform different functions.
		- Rendering is typically positioned downstream of simulation.
		- Rendering primarily consumes simulation state.
		- Rendering traditionally does not participate in simulation processing.

		Presentation Bullets:
		- Collision detection computes simulation behavior.
		- Rendering visualizes simulation results.
		- Rendering is typically downstream of simulation.
		- Rendering acts primarily as a consumer of simulation state.



	Paragraph 2

		Purpose:
		- Explain how GPU adoption affected rendering and collision detection.
		- Show that GPU execution did not fundamentally change their relationship.
		- Reinforce the separation between rendering and collision detection.
		- Identify the underutilization of rendering resources for collision detection.

		Contains:
		- Rendering and collision detection increasingly leverage GPU parallelism.
		- GPU-resident implementations became common.
		- Rendering and collision detection generally remain logically separated.
		- Collision detection executes in dedicated computational stages.
		- Rendering continues to generate visual output from simulation data.
		- Rendering contributes little or no information to collision detection.
		- GPU adoption alone does not integrate rendering and collision detection.

		Presentation Bullets:
		- Both rendering and collision detection moved to the GPU.
		- GPU residency did not eliminate their separation.
		- Collision detection uses dedicated compute stages.
		- Rendering remains focused on visualization.
		- Rendering contributes little to collision detection.




	Paragraph 3

		Purpose:
		- Review prior attempts to integrate rendering and simulation.
		- Acknowledge graphics-assisted approaches.
		- Explain the remaining limitation of existing methods.
		- Identify the architectural gap that motivates gPCD.

		Contains:
		- Prior work has explored rendering-simulation integration.
		- Approaches include:
			- Image-space techniques
			- GPU data structures
			- Graphics-assisted computation
		- Existing approaches still treat occupancy generation as a dedicated
		  collision-detection task.
		- Candidate construction remains independent of rasterization.
		- Rendering remains primarily downstream of simulation.
		- Rendering rarely acts as an active participant in collision detection.
		- An architectural opportunity remains unexploited.

		Presentation Bullets:
		- Prior work explored graphics-assisted simulation.
		- Image-space and GPU-assisted approaches exist.
		- Occupancy generation remains separate from rasterization.
		- Rendering is still primarily downstream of simulation.
		- An architectural gap remains.



	Paragraph 4

		Purpose:
		- Formulate the central architectural question.
		- Transition from prior work to the gPCD motivation.
		- Introduce graphics-pipeline occupancy generation as a possibility.
		- Establish concurrency between rendering and occupancy generation as a design objective.

		Contains:
		- Existing architectures separate occupancy generation from rendering.
		- This separation suggests an unexplored architectural opportunity.
		- Broad-phase occupancy generation may be integrated into the graphics pipeline.
		- Occupancy construction and rendering may proceed concurrently.
		- Graphics-pipeline resources may contribute directly to collision detection.

		Presentation Bullets:
		- Can occupancy generation be integrated into the graphics pipeline?
		- Can rendering and occupancy construction execute concurrently?
		- Can graphics resources participate directly in collision detection?
		- This question motivates gPCD.


	2.4 Gap in Existing Methods
		What architectural capability is missing?

	   Paragraph 1

		   Purpose:
		   - Summarize the findings of the previous background sections.
		   - Restate the traditional relationship between rendering and collision detection.
		   - Identify occupancy generation as a dedicated operation.
		   - Prepare the formal statement of the architectural gap.

		   Contains:
		   - Rendering and collision detection have traditionally remained separate.
		   - This separation persists in GPU-resident implementations.
		   - Rendering primarily consumes simulation state.
		   - Broad-phase occupancy generation is performed separately.
		   - Occupancy generation is treated as a dedicated collision-detection task.
		   - Existing architectures do not integrate occupancy generation into rendering.

		   Presentation Bullets:
		   - Rendering and collision detection remain separate.
		   - This separation persists on modern GPUs.
		   - Rendering consumes simulation state.
		   - Occupancy generation remains a dedicated collision-detection operation.



	Paragraph 2

		Purpose:
		- State the architectural gap explicitly.
		- Define the missing capability in existing methods.
		- Introduce the gPCD solution.
		- Connect the Background section to the Methodology section.

		Contains:
		- Existing methods do not integrate broad-phase occupancy generation directly into the graphics pipeline.
		- This absence constitutes the architectural gap.
		- gPCD addresses the gap through graphics-pipeline occupancy generation.
		- Occupancy construction and rendering proceed concurrently.
		- Graphics-pipeline execution produces the potentially colliding set (PCS).
		- The PCS is used during narrow-phase evaluation.
		- gPCD integrates rendering and occupancy generation into a unified architecture.

		Presentation Bullets:
		- Missing capability:
			- Graphics-pipeline occupancy generation.
		- gPCD integrates occupancy generation directly into the graphics pipeline.
		- Occupancy construction and rendering execute concurrently.
		- Graphics-pipeline execution produces the PCS.
		- gPCD closes the architectural gap.



	Paragraph 3

		Purpose:
		- Clarify the nature of the contribution.
		- Distinguish architectural novelty from algorithmic novelty.
		- Prevent misinterpretation of the work.
		- Reinforce the execution-architecture focus of gPCD.

		Contains:
		- gPCD does not introduce a new occupancy structure.
		- gPCD does not introduce new asymptotic complexity results.
		- Existing collision-detection techniques are retained.
		- The contribution lies in architectural organization.
		- Existing techniques are integrated into a unified GPU-driven framework.
		- The novelty is execution architecture rather than methodology.

		Presentation Bullets:
		- Not a new data structure.
		- Not a new complexity improvement.
		- Existing techniques are retained.
		- Novelty lies in execution architecture.
		- Unified GPU-driven framework.


3. Methodology
   How does gPCD work?

   3.1 Framework Overview
   Overall architecture.

   3.2 Graphics-Pipeline Occupancy Generation
   Broad-phase construction using the graphics pipeline.

   3.3 Corner-Based Occupancy Construction
   Corner generation and occupancy insertion.

   3.4 Duplicate Corner Suppression
   Preventing multiple insertions of the same particle
   into a cell occupancy list.

   3.5 Cell Construction and Occupancy Capacity
   Occupancy-list sizing from particle geometry.

   3.6 Cell Addressing and Flattening
   Cell indexing and memory access.

   3.7 PCS Construction
   Formation of the potentially colliding set.

   3.8 Source-Owned Narrow-Phase Evaluation
   Collision evaluation suitable for massively
   parallel execution.

4. Complexity Analysis
   How should the framework scale?

   4.1 Occupancy Generation
   Broad-phase complexity.

   4.2 PCS Construction
   Cost of candidate generation.

   4.3 Narrow-Phase Evaluation
   Collision-processing complexity.

   4.4 Storage Requirements
   Memory requirements and occupancy-list sizing.

   4.5 Overall Complexity
   Combined computational and storage complexity.

5. Experimental Setup
   What was tested?

   5.1 Hardware and Software
   Evaluation platform.

   5.2 Benchmark Configuration
   Test conditions and parameters.

   5.3 Benchmark Datasets
   Particle counts and benchmark-density studies.

   5.4 Performance Metrics
   Throughput, time-per-particle, and ARR metrics.

6. Results
   What happened?

   6.1 Throughput
   Measured particle throughput.

   6.2 Time per Particle
   Saturation behavior and scaling efficiency.

   6.3 Scaling Behavior
   Regression analysis and scaling trends.

   6.4 Locality Analysis
   Random versus sequential occupancy ordering.

   6.5 Density Studies
   Representative particle densities and achievable
   simulation scales.

   6.6 Comparison to Prior Work
   Relative performance and throughput.

7. Discussion
   Why did it happen?

   7.1 Interpretation of Scaling
   Meaning of measured scaling behavior.

   7.2 Saturation Behavior
   Utilization, saturation, and throughput-limited
   operating regions.

   7.3 Memory Locality
   Impact of occupancy ordering on performance.

   7.4 Implications for Large-Scale Particle Simulation
   Significance for engineering, scientific discovery,
   and future particle-based applications.

8. Conclusions
   What should the reader remember?

   Contributions.
   Key performance results.
   Practical implications.
   Future work.
   \end{verbatim}

\subsection{Purpose}

The purpose of this document is to:

\begin{itemize}
\item Maintain consistency across all manuscripts.
\item Prevent reintroduction of retired concepts.
\item Ensure technical claims remain supported by evidence.
\item Standardize notation and terminology.
\item Distinguish established results from future work.
\item Preserve continuity as the research evolves.
\end{itemize}

\subsection{Primary Contribution}

The principal architectural features of gPCD are:

\begin{itemize}
\item Graphics-pipeline generation of broad-phase occupancy information.
\item Duplicate-aware cell construction that prevents redundant particle occupancy entries.
\item Analytically bounded occupancy-list sizing derived from particle geometry.
\item Source-owned narrow-phase collision evaluation suitable for massively parallel execution.
\item Demonstrated scalability to more than ten million particles on a single CPU-GPU system.
\end{itemize}


\subsection{Abstract Requirements}

The abstract shall be written after completion of the manuscript
review process.

The abstract shall answer the following questions:

\begin{enumerate}
\item What problem is addressed?
\item What is the primary contribution?
\item How does the method operate?
\item What results were demonstrated?
\item Why are the results significant?
\end{enumerate}

The abstract shall not introduce terminology, concepts, or claims
that are not subsequently developed and supported within the manuscript.

\subsection{Paper Scope}

The current manuscript is exclusively a particle collision-detection
paper describing the gPCD framework.

Topics that belong within the scope of this manuscript include:

\begin{itemize}
\item Broad-phase collision detection.
\item Cell-array occupancy generation.
\item Graphics-pipeline occupancy generation.
\item Duplicate elimination.
\item Narrow-phase collision evaluation.
\item Computational complexity.
\item GPU implementation.
\item Performance and scalability.
\end{itemize}

Topics that are explicitly outside the scope of this manuscript include:

\begin{itemize}
\item Field Particle Method dynamics.
\item Internal momentum representation.
\item Thermodynamic modeling.
\item Pressure and temperature formulation.
\item Phase-change modeling.
\item Converging-diverging nozzle analysis.
\item Choking behavior.
\item Flow diagnostics.
\end{itemize}

Such topics may be referenced only when necessary to provide context and
shall not be presented as primary contributions of the manuscript.

\subsection{Section Responsibilities}

Each major section of the manuscript is intended to answer a distinct
question.

\begin{itemize}
\item Methodology: What does the algorithm do?
\item Complexity Analysis: How should the algorithm scale?
\item Experimental Setup: How was the algorithm evaluated?
\item Results: What performance was observed?
\item Discussion: Why was this behavior observed and what are its implications?
\item Conclusions: What should the reader take away from the work?
\end{itemize}

To avoid redundancy, complexity analysis, performance measurements, and
interpretation are separated into their respective sections whenever
possible.

\subsection{Revision Procedure}

Before major modifications are made to any manuscript:

\begin{enumerate}
\item Proposed terminology changes shall be reviewed against this document.
\item New technical claims shall be added here before being incorporated into manuscripts.
\item Retired concepts shall be recorded in Section~\ref{sec:retired_concepts}.
\item Publication-specific deviations shall be documented and justified.
\end{enumerate}

\section{Figures}

Figures are classified according to their role within the manuscript.
Each figure type is intended to support a specific section of the paper
and should be placed where it contributes most directly to the reader's
understanding.

\begin{table}[h!]
\centering
\caption{Figure classification and recommended placement.}
\label{tab:figure_classification}
\begin{tabular}{p{3.0cm} p{6.5cm} p{3.0cm}}
\hline
\textbf{Figure Type} & \textbf{Purpose} & \textbf{Typical Section} \\
\hline

Architecture Figure &
Describes system organization, workflow, or data flow. &
Methodology \\

Benchmark-Dataset Figure & Illustrates benchmark configurations, datasets, particle distributions,
or test environments used during evaluation. & Experimental Setup \\

Performance Figure &
Presents measured results such as timing, throughput, scaling,
efficiency, or resource utilization. &
Results \\

Interpretation Figure & Provides explanation of observed behavior, comparison of trends, or
discussion of performance characteristics. & Discussion \\

\hline
\end{tabular}
\end{table}

The purpose of this classification is to maintain a clear distinction
between describing the framework, defining the evaluation, presenting
measurements, and interpreting results. Figures should generally appear
in the section corresponding to their primary role within the paper.

Examples within this work include architectural overview diagrams in
the Methodology section, benchmark-density visualizations in
Experimental Setup, throughput and timing plots in Results, and
saturation or scaling interpretations in Discussion.

\begin{table}[h!]
\centering
\caption{Current figure inventory and recommended placement within the manuscript.}
\label{tab:figure_inventory}
\begin{tabular}{p{4.0cm} p{3.5cm} p{3.5cm}}
\hline
\textbf{Figure} & \textbf{Type} & \textbf{Section} \\
\hline

OverviewArchitecture &
Architecture Figure &
Methodology \\

Corner Occupancy / Duplicate Collision Figure &
Methodology Figure &
Methodology \\

Cell Flattening Figure (if added) &
Methodology Figure &
Methodology \\

Benchmark Particle Density Figure &
Benchmark/Dataset Figure &
Experimental Setup \\

Throughput Plots &
Performance Figure &
Results \\

Time-per-Particle Plots &
Performance Figure &
Results \\

Scaling Fits &
Performance Figure &
Results \\

Saturation Explanation Figure (if added) &
Interpretation Figure &
Discussion \\

\hline
\end{tabular}
\end{table}

\section{Performance Regions}

Throughout this work, throughput curves are interpreted in terms of
three performance regions:

\begin{enumerate}

\item \textbf{Utilization Region}

The utilization region corresponds to low particle counts where GPU
resources are not yet fully utilized. In this region, throughput
increases rapidly as particle count increases because additional work
improves parallel occupancy and hardware utilization.

\item \textbf{Peak Throughput Region}

The peak throughput region corresponds to the neighborhood surrounding
the maximum measured throughput. Within this region, the framework
achieves its highest particle-processing rate and transitions from
utilization-limited behavior toward throughput-limited behavior.

\item \textbf{Throughput-Limited Region}

The throughput-limited region corresponds to particle counts beyond the
peak throughput region. In this region, additional particles no longer
increase throughput and may produce a gradual decline in throughput as
execution costs continue to grow.

\end{enumerate}

Observed transition locations are benchmark dependent and are
identified directly from measured throughput curves rather than imposed
a priori.

\subsection{Sequential vs Random Percent increase}

The percent increase comparison is calculated by \ref{pci}
\begin{equation}
\label{pci}
P_{\mathrm{locality}}
=
\frac{t_{\mathrm{random}}-t_{\mathrm{sequential}}}
     {t_{\mathrm{sequential}}}
\times 100\%.
\end{equation}

\newpage
\section{Comparative Analysis}

\section{Comparative Analysis}

The time-per-particle (TPP) values reported in
Table~\ref{tab:tpp_comparison} were derived from performance metrics
reported in the corresponding publications. Because the surveyed works
span multiple decades, hardware platforms, benchmark configurations,
and particle-processing architectures, the resulting values should be
interpreted as approximate indicators of reported particle-processing
efficiency rather than direct performance rankings.

For each reference, the reported particle count (N), execution time
(t), frame rate, or throughput metric was extracted from the
publication and converted to a common time-per-particle metric using

\begin{equation}\mathrm{TPP}=\frac{t}{N},\end{equation}

where (t) is expressed in nanoseconds and (N) is the number of
particles processed.

The following subsections document the source data and calculations
used to construct Table~\ref{tab:tpp_comparison}.

\subsection{Latta (2004)}

Reference:
\cite{latta}

Reported Metrics:

\begin{itemize}
\item Particle count:
\item Reported execution time:
\item Reported frame rate:
\item Hardware:
\end{itemize}

Calculation:

\begin{equation}
\mathrm{TPP}_{Latta}=33.33;\mathrm{ns/particle}.
\end{equation}

Comments:

\begin{itemize}
\item
\end{itemize}

\subsection{Kipfer (2004)}

Reference:
\cite{kipfer}

Reported Metrics:

\begin{itemize}
\item Particle count:
\item Reported execution time:
\item Reported frame rate:
\item Hardware:
\end{itemize}

Calculation:

\begin{equation}
\mathrm{TPP}_{Kipfer}=390.00;\mathrm{ns/particle}.
\end{equation}

Comments:

\begin{itemize}
\item
\end{itemize}

\subsection{Zhang (2008)}

Reference:
\cite{zhang}

Reported Metrics:

\begin{itemize}
\item Particle count:
\item Reported execution time:
\item Reported frame rate:
\item Hardware:
\end{itemize}

Calculation:

\begin{equation}
\mathrm{TPP}_{Zhang}=\mathrm{ns/particle}.
\end{equation}

\subsection{Anderson (2008)}

Reference:
\cite{anderson}

Reported Metrics:

\begin{itemize}
\item Particle count:
\item Reported execution time:
\item Reported frame rate:
\item Hardware:
\end{itemize}

Calculation:

\begin{equation}
\mathrm{TPP}_{Anderson}=1400.00;\mathrm{ns/particle}.
\end{equation}

\subsection{Liu (2010)}

Reference:
\cite{liu}

Reported Metrics:

\begin{itemize}
\item Particle count:
\item Reported execution time:
\item Reported frame rate:
\item Hardware:
\end{itemize}

Calculation:

\begin{equation}
\mathrm{TPP}_{Liu}=167.70;\mathrm{ns/particle}.
\end{equation}

\subsection{Govender (2015)}

Reference:
\cite{govender}

Reported Metrics:

\begin{itemize}
\item Particle count:
\item Reported execution time:
\item Reported frame rate:
\item Hardware:
\end{itemize}

Calculation:

\begin{equation}
\mathrm{TPP}_{Govender}=943.00;\mathrm{ns/particle}.
\end{equation}

\subsection{Joselli (2015)}

Reference:
\cite{joselli}

Reported Metrics:

\begin{itemize}
\item Particle count:
\item Reported execution time:
\item Reported frame rate:
\item Hardware:
\end{itemize}

Calculation:

\begin{equation}
\mathrm{TPP}_{Joselli}=\mathrm{ns/particle}.
\end{equation}

\subsection{Franklin (2017)}

Reference:
\cite{franklin}

Reported Metrics:

\begin{itemize}
\item Particle count:
\item Reported execution time:
\item Reported frame rate:
\item Hardware:
\end{itemize}

Calculation:

\begin{equation}
\mathrm{TPP}_{Franklin}=33.00;\mathrm{ns/particle}.
\end{equation}

\subsection{Mittal (2020)}

Reference:
\cite{mittal}

Reported Metrics:

\begin{itemize}
\item Particle count:
\item Reported execution time:
\item Reported frame rate:
\item Hardware:
\end{itemize}

Calculation:

\begin{equation}
\mathrm{TPP}_{Mittal}=16.50;\mathrm{ns/particle}.
\end{equation}

\subsection{Bell (2025)}

Reported Metrics:

\begin{itemize}
\item Particle count:
\numprint{11088896}
\item Total execution time:
84.47 ms
\item Hardware:
RTX 3500 Ada Generation Laptop GPU
\end{itemize}

Calculation:

\begin{equation}
\mathrm{TPP}_{Bell}=\frac{84.47\times10^{6}}{\numprint{11088896}}=7.62;\mathrm{ns/particle}.
\end{equation}

Comments:

\begin{itemize}
\item Includes occupancy generation, occupancy-list traversal,
narrow-phase evaluation, and rendering.
\item Reported using measured execution time from the complete
framework.
\end{itemize}



\subsection{Uniform Grids}
Q: Uniform grids with direct insertion are also O(N).
A: The contribution is not a new asymptotic complexity class. The contribution is graphics-pipeline occupancy generation, duplicate-aware occupancy construction, and a GPU-driven architecture that achieves O(N) broad-phase generation while maintaining favorable memory locality."

\section{Audience and Explanation Level}

The target audience for this manuscript consists primarily of ACM
Transactions on Graphics (TOG) reviewers, computer graphics
researchers, GPU-computing researchers, and collision-detection
specialists.

Assume familiarity with standard graphics and GPU-computing concepts,
including but not limited to:

\begin{itemize}
\item graphics pipelines,
\item vertex and fragment shaders,
\item compute shaders,
\item shader storage buffer objects (SSBOs),
\item GPU residency,
\item broad-phase collision detection,
\item narrow-phase collision detection, and
\item common spatial data structures.
\end{itemize}

Consequently, methodology sections should focus on explaining the
novel contributions of gPCD rather than providing tutorial-level
descriptions of established graphics concepts.

Writing Rule:

\begin{quote}
Explain novelty. Reference common knowledge.
\end{quote}

When drafting methodology sections, ask:

\begin{enumerate}
\item Is this sentence explaining graphics concepts already known to the
target audience?
\item Or is this sentence explaining what is unique about gPCD?
\end{enumerate}

If the material can be found in a standard graphics or GPU-programming
textbook, it likely does not belong in the methodology. If the material
describes a design decision, organizational principle, or algorithmic
component unique to gPCD, it likely does belong in the methodology.

\section{Drafting Note}

The conventions contained in this document are maintained as part of the
working manuscript during the rough-draft and revision stages. Their
purpose is to preserve notation consistency, terminology decisions,
architectural interpretations, and reviewer-driven revisions throughout
development.

Items contained in this document are not necessarily intended for
publication. As the manuscript matures, notation and terminology should
be introduced at the point of first use within the manuscript, and
working notes should be removed from the final submission.

Consequently, this document serves as both a notation reference and a
development history for the gPCD manuscript.

%================================================================================
\section{Document Conventions and Research Governance}

This document serves as the authoritative reference for terminology,
notation, technical claims, algorithm descriptions, and publication
standards used throughout the gPCD and Field Particle Method research
program. When conflicts arise between a manuscript and this document,
this document shall take precedence until formally revised.


%================================================================================
\section{Approved Terminology}

\subsection{Terminology Convention: Potentially Colliding Set (PCS)}

Within collision-detection literature, PCS denotes the
\emph{Potentially Colliding Set}. Within the gPCD framework, the
populated cell array constitutes the PCS after occupancy generation has
been completed. Consequently, the manuscript should distinguish between
the cell array as a data structure and the PCS as the broad-phase result
represented by that populated structure.

When discussing implementation details, prefer the terms
\emph{cell array} and \emph{occupancy list}. The term
\emph{Potentially Colliding Set (PCS)} should be reserved for
discussions of broad-phase candidate generation and collision-detection
theory.

\subsection*{PCS Convention}

Within the gPCD framework, occupancy generation directly constructs the
potentially colliding set (PCS). The populated cell array therefore
constitutes the PCS and no separate candidate-construction stage is
performed between broad-phase occupancy generation and narrow-phase
evaluation.

\begin{equation}
\text{Broad Phase}
\rightarrow
\text{Candidate Construction}
\rightarrow
\text{PCS}.
\end{equation}

Within the gPCD framework, occupancy generation directly constructs the
potentially colliding set,

\begin{equation}
\text{Occupancy Generation}
\rightarrow
\text{PCS}.
\end{equation}

Consequently, the populated cell array constitutes the PCS and no
separate candidate-construction stage is performed prior to
narrow-phase evaluation.

\subsection{Symbols}
\begin{tabular}{|l|c|}
\hline
Topic & Status \\
\hline
$P$ & number of particles\\
$S$ & side length of cubic cell array\\
$C$ & total number of cells = $S^3$\\
$r$ & particle radius\\
$(i,j,k)$ & floating-point cell coordinates\\
$(X,Y,Z)$ & integer cell coordinates\\
$I$ & flattened cell index\\
\hline
\end{tabular}

\subsection{Symbol Conventions}

The gPCD framework utilizes three coordinate systems and one indexing
system. Distinguishing between these representations is essential for
understanding occupancy construction and complexity analysis.

\subsubsection{Physical Particle Coordinates}

Particle positions are represented in continuous physical space using
floating-point coordinates,

\begin{equation}
(x,y,z).
\end{equation}

These coordinates correspond to the particle center location and are
used for geometric calculations.

\subsubsection{Cell Coordinates}

Cell coordinates are obtained by rounding physical coordinates to the
nearest integer cell location,

\begin{equation}
(X,Y,Z).
\end{equation}

The cell coordinates are stored as unsigned integers and identify the
cell occupied by a particle corner.

Valid cell coordinates satisfy

\begin{equation}
X \geq 1,
\qquad
Y \geq 1,
\qquad
Z \geq 1.
\end{equation}

A coordinate value of zero is reserved as an invalid location.
\begin{equation}
X = 0 \;\lor\; Y = 0 \;\lor\; Z = 0.
\end{equation}

indicates that the coordinate does not reference a valid cell.

\subsubsection{Flattened Cell Index}

The integer cell coordinates are converted to a one-dimensional
flattened index,

\begin{equation}
I.
\end{equation}

The flattened index uniquely identifies a cell within the cell array
and is stored as an unsigned integer.

The value

\begin{equation}
I = 0
\end{equation}

is reserved as an invalid index.

\subsection{Cell Coordinate Convention}

Cell coordinates are restricted to positive integer values. The valid
cell domain therefore begins at $(1,1,1)$ rather than $(0,0,0)$.
Coordinates containing a zero-valued component are considered invalid.
This convention prevents corner evaluations associated with particles
located adjacent to domain boundaries from producing negative cell
coordinates during occupancy construction.

\subsubsection{Cell Array Dimensions}

The dimensions of the cell array are represented by

\begin{equation}
(S_x,S_y,S_z),
\end{equation}

where each quantity denotes the number of cells in the corresponding
coordinate direction.

The total number of cells is therefore

\begin{equation}
C = S_x S_y S_z.
\end{equation}

For cubic cell arrays,

\begin{equation}
S_x = S_y = S_z = S,
\end{equation}

and the total cell count becomes

\begin{equation}
C = S^3.
\end{equation}

Consequently, cell-array memory requirements scale as

\begin{equation}
\mathcal{O}(S^3).
\end{equation}

\subsubsection{Storage Conventions}

The value zero is reserved as a sentinel value throughout the gPCD
framework.

\begin{itemize}
\item A cell coordinate value of zero indicates an invalid cell location.
\item A flattened index of zero indicates an invalid cell index.
\item An occupancy-list entry of zero indicates an empty array location.
\end{itemize}

As a result, valid cells and valid occupancy entries begin at one.

\subsubsection{Coordinate Transformation Pipeline}

Occupancy construction follows the transformation sequence

\begin{equation}
(x,y,z)
;\rightarrow;
(X,Y,Z)
;\rightarrow;
I,
\end{equation}

where floating-point particle coordinates are rounded to integer cell
coordinates and subsequently converted into flattened cell indices for
occupancy insertion.



\subsection{Bounding Volume Terminology}

The term AABB (Axis-Aligned Bounding Box) shall be used throughout all
manuscripts. Any occurrence of AARB shall be treated as a typographical
error and corrected to AABB.
%================================================================================
\section{Document Structure}

\subsection{Architecture Hierarchy}
\begin{enumerate}
    \item Graphics-Pipeline Occupancy Generation
    \item Vertex-Stage Occupancy Construction
    \item Corner-Based Occupancy Generation
    \item Cell Addressing and Flattening
    \item Potentially Colliding Set Construction
    \item Narrow-Phase Evaluation
\end{enumerate}

\subsection{Numerical Placeholders}

During manuscript development, unresolved numerical values should be
represented using
\begin{verbatim}
    \nc{Field Name}
\end{verbatim}
rather than hard-coded values.

Examples:
\begin{verbatim}
    \nc{Maximum Particle Count}
    \nc{Peak Throughput}
    \nc{Scaling Exponent}
    \nc{Time per Particle}
\end{verbatim}
Placeholders allow prose development to proceed independently of
numerical refinement and provide a visible checklist of values that
must be finalized before submission.
%================================================================================
\section{Collision-Detection Architecture}

The manuscript shall distinguish between collision-detection algorithms
and collision-detection architectures.

A collision-detection framework consists of broad-phase and narrow-phase
stages. The broad phase generates collision candidates, while the
narrow phase evaluates those candidates and determines whether a
collision exists.

\subsection*{Cell-Array Architecture}

The gPCD cell array is organized as a fixed-size occupancy structure.
Let $S$ denote the side length of the cubic cell array and let $W$
denote the occupancy-list width assigned to each cell. The cell array
therefore contains

\begin{equation}
S^3
\end{equation}

cells, each capable of storing up to $W$ occupancy entries.

Consequently, the cell array may be viewed as a two-dimensional storage
structure consisting of a cell dimension of length $S^3$ and an
occupancy dimension of width $W$,

\begin{equation}
\text{Cell Array} \equiv [S^3][W].
\end{equation}

The parameter $W$ is a fixed implementation constant selected to
accommodate the maximum expected occupancy within a cell.

\subsection{Traditional Architectures}

Traditional collision-detection systems typically perform broad-phase
and narrow-phase operations independently of the rendering pipeline.
Collision-detection results are subsequently transferred to the
rendering system for visualization.

\begin{verbatim}
Simulation Data
|
+-- Broad Phase
|
+-- Narrow Phase
|
+-- Rendering Pipeline
|
+-- Visual Output
\end{verbatim}

Rendering functions primarily as a consumer of simulation state and
does not participate directly in collision detection.

\subsection{GPU-Resident Architectures}

GPU-resident approaches perform broad-phase and narrow-phase operations
on the GPU while maintaining a separate rendering pipeline.

\begin{verbatim}
Simulation Data
|
+-- GPU Broad Phase
|
+-- GPU Narrow Phase
|
+-- GPU Rendering Pipeline
|
+-- Visual Output
\end{verbatim}

Although collision detection and rendering execute on the same device,
they remain logically distinct processes.

\subsection{gPCD Architecture}

The gPCD framework integrates occupancy generation directly into the
graphics pipeline. The vertex shader participates in broad-phase collision
detection while simultaneously supporting particle rendering.

\begin{verbatim}
Simulation Data
|
+-- Graphics Pipeline
|       |
|       +-- Point-Particle Rendering
|       |
|       +-- Vertex-Based Occupancy Generation
|               |
|               +-- Cell-Array Occupancy
|               |
|               +-- Collision Candidates
|
+-- Compute Pipeline
|
+-- Narrow-Phase Collision Evaluation
|
+-- Collision Results
\end{verbatim}

Within gPCD, rendering is not merely a consumer of collision-detection
results. Instead, the graphics pipeline forms part of the broad-phase
collision-detection process.

Collision-Detection Data-Flow Architectures
\lstset{
  language=TeX,
  basicstyle=\ttfamily,
}

\begin{lstlisting}

1. CPU-Resident / CPU-Driven Pipeline

Simulation Data
    |
    +-- CPU Broad Phase
    |
    +-- CPU Narrow Phase
    |
    +-- Particle/Geometry Transfer
    |
    +-- GPU Rendering Pipeline
            |
            +-- Visual Output


2. GPU-Assisted or GPU-Resident Pipeline

Simulation Data
    |
    +-- GPU Broad Phase
    |
    +-- GPU Narrow Phase
    |
    +-- GPU Rendering Pipeline
            |
            +-- Visual Output


3. gPCD Pipeline

Simulation Data
    |
    +-- Graphics Pipeline
    |       |
    |       +-- Point-Particle Rendering
    |       |
    |       +-- Vertex-Based Occupancy Generation
    |               |
    |               +-- Cell-Array Occupancy
    |               |
    |               +-- Collision Candidates
    |
    +-- Compute Pipeline
            |
            +-- Narrow-Phase Collision Evaluation
            |
            +-- Collision Results

\subsection{gPCD}
\end{lstlisting}


\subsection{Architectural Classification}

The primary organizational contribution of gPCD is architectural rather
than structural. The novelty does not arise from the cell array itself, but from
integrating occupancy generation into the graphics pipeline such that
occupancy construction is performed within the vertex stage during
broad-phase collision detection.



\subsection{Graphics Pipeline Execution}

The graphics pipeline shall be described as performing
concurrent point-particle rendering and broad-phase occupancy
generation. The term ``hierarchical parallelism'' shall not
be used to describe gPCD execution.

Vertex and fragment invocations are treated as independent
stages of the graphics pipeline and not as a nested or
parent-child execution hierarchy.



The term \textbf{gPCD} refers to the GPU-driven particle collision detection
system consisting of:

\begin{itemize}
\item Graphics-pipeline broad-phase occupancy generation.
\item Cell-array collision candidate construction.
\item Compute-shader narrow-phase collision evaluation.
\item Fully GPU-resident execution.
\end{itemize}

\subsection{Broad Phase}

The broad phase refers exclusively to occupancy generation and collision
candidate identification using the cell-array structure.

\subsection{Narrow Phase}

The narrow phase refers exclusively to exact collision testing and collision
resolution performed within compute shaders.

\subsection{Cell Array}

The term \textbf{cell array} shall be used instead of spatial grid whenever
referring to the gPCD occupancy structure.

\subsection*{Cell-Array Architecture}

The gPCD cell array consists of a fixed number of cells and a fixed
occupancy-list width per cell. Let $S$ denote the side length of the
cubic cell array and let $W$ denote the occupancy-list width assigned
to each cell.

The total number of cells is

\begin{equation}
S^3.
\end{equation}

Each cell stores up to $W$ occupancy entries, yielding a conceptual
storage organization of

\begin{equation}
\text{Cell Array} \equiv [S^3][W].
\end{equation}

The parameter $W$ is a fixed implementation constant selected to
accommodate the maximum expected occupancy within a cell.

\section{Related Work Objectives}

The Background and Related Work section shall establish the context
necessary to evaluate the novelty and performance of gPCD.

The section shall:

\begin{enumerate}
\item Introduce broad-phase collision-detection methodologies.
\item Summarize GPU-based collision-detection approaches.
\item Identify limitations or distinctions relevant to gPCD.
\item Explain how gPCD differs from existing approaches.
\end{enumerate}

The section shall not function as a comprehensive literature survey.
Its purpose is to provide sufficient context for evaluating the
organizational and algorithmic contributions of the framework.



%================================================================================
\section{Experimental Claims}

The Experimental Results section shall validate the effectiveness of the
gPCD framework rather than establish novelty.

Experimental evidence shall support the following claims:

\begin{enumerate}
\item The framework processes large particle counts efficiently.
\item The framework reaches a stable throughput-limited operating regime.
\item The framework exhibits near-linear scaling behavior.
\end{enumerate}

Throughput, time-per-particle, curvature analysis, and log-log
regression shall be treated as complementary measurements supporting
these claims rather than as independent studies.

Performance results shall be presented as evidence supporting the
organizational and algorithmic contributions of the framework.


%================================================================================
\section{Complexity Claims}

The gPCD framework exhibits distinct access, computational, and storage
complexities. These quantities shall be treated independently throughout
the manuscript.

The complexity analysis assumes a fixed maximum particle size
consistent with the eight-cell occupancy constraint.

\subsection{Access Complexity}

Cell lookup is performed through direct conversion of integer cell
coordinates into a flattened cell index. No search, traversal, or
hierarchical lookup is required. Consequently, cell access complexity is

\begin{equation}
\mathcal{O}(1).
\end{equation}

Similarly, each particle corner is processed independently using a fixed
sequence of coordinate rounding and index generation operations.
Therefore, corner evaluation complexity is

\begin{equation}
\mathcal{O}(1).
\end{equation}

\subsection{Computational Complexity}

Each particle contributes a fixed number of corner evaluations. Since
the gPCD framework evaluates exactly eight corners per particle, the
cost of occupancy generation scales linearly with the number of
particles (P),

\begin{equation}
\mathcal{O}(8P)=\mathcal{O}(P).
\end{equation}

Because broad-phase occupancy construction is composed entirely of these
fixed corner evaluations, broad-phase construction complexity is also

\begin{equation}
\mathcal{O}(P).
\end{equation}

\subsection{Storage Complexity}

Storage complexity is determined by the dimensions of the cell array and
is independent of particle count.

For a cell array with dimensions

\begin{equation}
(S_x,S_y,S_z),
\end{equation}

the total number of cells is

\begin{equation}
C=S_xS_yS_z.
\end{equation}

Therefore, cell-array memory requirements scale as

\begin{equation}
\mathcal{O}(S_xS_yS_z).
\end{equation}

For cubic domains,

\begin{equation}
S_x=S_y=S_z=S,
\end{equation}

yielding

\begin{equation}
C=S^3,
\end{equation}

and consequently

\begin{equation}
\mathcal{O}(S^3).
\end{equation}

\subsection{Complexity Convention}

The manuscript shall distinguish between:

\begin{itemize}
\item \textbf{Access Complexity} -- the cost of retrieving or updating a
cell location.
\item \textbf{Computational Complexity} -- the cost of processing
particles and constructing occupancy information.
\item \textbf{Storage Complexity} -- the memory required to store the
cell array.
\end{itemize}

Statements regarding constant-time cell access shall not be interpreted
as statements regarding memory growth. The gPCD framework provides
constant-time cell access,

\begin{equation}
\mathcal{O}(1),
\end{equation}

while cell-array storage grows according to

\begin{equation}
\mathcal{O}(S^3)
\end{equation}

for cubic domains.

%================================================================================

\section{Novelty Statement}

The primary novelty of gPCD lies in the integration of occupancy
generation into the graphics pipeline for broad-phase collision
detection.The method constructs the occupancy structure as part of the graphics pipeline,
allowing point-particle rendering and collision-candidate generation
to proceed concurrently within a fully GPU-resident workflow.

The novelty of the manuscript shall not be attributed to cell arrays,
uniform grids, AABBs, rasterization, fragment shaders, compute shaders,
or broad-phase/narrow-phase decomposition individually, as these
concepts are established within the literature.

\subsection{Novelty Classification}

The manuscript shall distinguish between:

\begin{enumerate}
\item Established components,
\item Organizational contributions, and
\item Performance outcomes.
\end{enumerate}

Established components may contribute significantly to performance
without constituting the primary novelty of the manuscript.

The primary novelty of gPCD is expected to arise from the organization
and interaction of established components within a graphics-pipeline
occupancy-generation framework for GPU-resident collision detection.

Performance results shall be treated as validation of the organization
rather than as novelty by themselves.

\subsection{Primary Organizational Contribution}

The primary organizational contribution of gPCD is the integration of
broad-phase occupancy generation into the graphics pipeline itself.
Occupancy construction is performed within the vertex stage, where
particle geometry is used to generate the occupancy information required
for broad-phase collision detection.

Because occupancy generation is integrated into graphics-pipeline
execution, point-particle rendering and broad-phase construction proceed
concurrently on the GPU. This concurrency is considered a consequence of
the organizational design rather than a separate contribution.


The novelty of the method arises from the organization of established
graphics-pipeline components into a collision-detection workflow rather
than from the individual components themselves.
\newpage
\lstset{
  language=TeX,
  basicstyle=\ttfamily,
}

\begin{lstlisting}

Primary Organizational Contribution
|
+-- Graphics pipeline integrated into collision detection
	|
	+-- Vertex generates occupancy information
    |
    +-- Rendering becomes part of the algorithm
    |
    +-- Supporting Algorithmic Contribution
        |
        +-- Corner-based occupancy construction
            |
            +-- Corner rounding
            +-- Integer cell coordinates
            +-- Flattened indexing
            |
            +-- Emergent Properties
                |
                +-- Fixed eight-cell occupancy guarantee
                +-- Bounded occupancy generation
                +-- Predictable workload

Methodological Requirements
|
+-- Atomics
+-- Duplicate elimination
+-- Memory structure

Performance Outcomes
|
+-- Throughput
+-- Near-linear scaling
+-- GPU residency

\end{lstlisting}
\subsection{Candidate Supporting Contributions}

\begin{itemize}
\item Fixed eight-cell occupancy constraint.
\item Corner-based occupancy construction through coordinate
      rounding and cell indexing.
\end{itemize}

gPCD constructs cell occupancy through a corner-based formulation in
which particle corner coordinates are rounded to integer cell
coordinates and converted to flattened cell indices. Occupancy
generation is therefore reduced to a fixed set of corner evaluations
rather than cell-overlap traversal or neighborhood discovery.

The resulting occupancy-construction procedure is simple, bounded, and
well suited to massively parallel execution. Although atomic operations
are required during occupancy insertion, the fixed-size formulation
provides predictable computational behavior and contributes to the
overall efficiency of the broad-phase algorithm.


\subsubsection{Supporting Contributions}


\begin{itemize}
\item Corner-based occupancy construction through corner-coordinate
      rounding and flattened indexing.
\item Fixed eight-cell occupancy guarantee arising from the occupancy
      construction formulation and particle-size constraint.
\end{itemize}

\subsection{Contribution Hierarchy}
\begin{lstlisting}
Organizational Contributions->Algorithmic Contributions
->Emergent Properties->Methodological Requirements
->Performance Outcomes
\end{lstlisting}

\subsection*{Contribution Clarification}

The primary organizational contribution of gPCD is the integration of
broad-phase occupancy generation into the graphics pipeline.

Occupancy construction is performed within the vertex stage of the
graphics pipeline through corner-based occupancy generation and cell
determination. The fragment stage remains responsible for rendering and
visualization.

Consequently, manuscript terminology should preferentially use
\emph{graphics-pipeline occupancy generation} rather than
\emph{rasterization-based occupancy generation}, unless a discussion
specifically concerns rasterization itself.

This distinction is important because the contribution of gPCD arises
from the integration of collision-detection operations into the graphics
pipeline, while occupancy construction is implemented within the vertex
stage rather than the fragment stage.
%================================================================================
\section{Methodological Requirements}

\subsection{Correctness Mechanisms}

Certain algorithmic components exist to ensure correctness of the
occupancy-generation process and shall not be treated as primary
novel contributions.

Examples include duplicate-elimination procedures required to
remove repeated occupancy and collision evaluations generated by
the occupancy-construction methodology.

\subsection{Memory}
%Place holder

%================================================================================
\section{Approved Technical Claims}

The following statements are currently approved for publication use.

\begin{itemize}
\item gPCD is fully GPU resident.
\item Broad-phase occupancy generation occurs during graphics-pipeline execution.
\item Point-particle rendering and broad-phase occupancy generation occur concurrently.
\item A particle may occupy at most eight cells.
\item Cell-array indexing provides constant-time cell access.
\item The system demonstrates near-linear scaling over the tested range.
\end{itemize}

%================================================================================
\section{Retired Concepts}
\label{sec:retired_concepts}

The following concepts are considered obsolete and shall not be reintroduced
without formal revision of this document.

\begin{itemize}
\item Hierarchical GPU execution.
\item Hierarchical broad-phase processing.
\item Hierarchical occupancy generation.
\item Nested graphics-pipeline execution.
\item Unsupported novelty claims.
\item Unverified overlap claims presented as measured results.
\end{itemize}

%================================================================================
\section{Writing Standards}

\begin{itemize}
\item Manuscripts shall follow ACM formatting conventions unless otherwise required.
\item Equations shall be numbered when referenced.
\item Variables shall be defined upon first use.
\item Technical claims shall be supported by either derivation, measurement, or citation.
\item Future work shall be clearly separated from demonstrated results.
\end{itemize}

%================================================================================
\section{Research Separation}

\subsection{gPCD Paper}

The gPCD paper shall focus on:

\begin{itemize}
\item Collision detection.
\item Graphics-pipeline occupancy generation.
\item Cell-array construction.
\item GPU performance.
\item Scalability.
\end{itemize}

\subsection{Field Particle Method Paper}

The Field Particle Method paper shall focus on:

\begin{itemize}
\item Geometric dynamics.
\item Internal momentum representation.
\item Thermo-fluid behavior.
\item Boundary-condition treatment.
\item Physical interpretation of emergent behavior.
\end{itemize}

\subsection{Active Manuscript Rule}

Only the current manuscript shall be used when evaluating terminology,
technical claims, notation, and structure. Archived files, legacy drafts,
and historical manuscripts shall not be used as evidence for current
terminology decisions.

Historical files may be consulted only for recovering lost content or
tracking the evolution of a concept.

\subsection{Item 1 Status}

\textbf{Resolved.}

The term ``hierarchical parallelism'' has been retired as a description
of gPCD. The graphics pipeline shall instead be described as performing
concurrent point-particle rendering and broad-phase occupancy generation.

Acceptable terminology includes:

\begin{itemize}
\item Graphics-pipeline parallelism
\item Concurrent rendering and occupancy generation
\item Graphics Pipline-based occupancy generation
\item GPU-resident broad-phase construction
\end{itemize}


\section{Manuscript Review Status}

\begin{tabular}{|l|c|}
\hline
Topic & Status \\
\hline
Hierarchical Parallelism  &     Resolved \\
AABB Terminology          &    Resolved \\
Paper Scope                &   Resolved\\
Novelty Structure          &   Resolved\\
Contribution Hierarchy     &   Resolved\\
Complexity Discussion  &       Resolved\\
Experimental Claims     &      Pending\\
Primary Contribution   &       Pending\\
Abstract                &      Pending\\
\hline
\end{tabular}

